\section{Cancellations in alternating sums}

\subsection{Acceptable sets, partners, and signs}

\begin{definition}
\label{def.acceptableSets}
\lean{AlgebraicCombinatorics.SignedCounting.acceptableSets}
\leanhelper
\leanok
The \emph{acceptable sets} for parameters $n, m \in \mathbb{N}$ are the
subsets of $\{0, 1, \ldots, n-1\}$ with cardinality at most $m$:
\[
\mathcal{A}(n, m) := \bigl\{I \subseteq \{0, \ldots, n-1\} : |I| \leq m\bigr\}.
\]
\end{definition}

\begin{definition}
\label{def.partner}
\lean{AlgebraicCombinatorics.SignedCounting.partner}
\leanhelper
\leanok
The \emph{partner} of a finite set $I$ of natural numbers is the
symmetric difference $I' := I \bigtriangleup \{0\}$.
If $0 \in I$, then $I' = I \setminus \{0\}$;
if $0 \notin I$, then $I' = I \cup \{0\}$.
\end{definition}

\begin{definition}
\label{def.setSign}
\lean{AlgebraicCombinatorics.SignedCounting.setSign}
\leanhelper
\leanok
The \emph{sign} of a finite set $I$ of natural numbers is
$\operatorname{sign}(I) := (-1)^{|I|}$.
\end{definition}

\begin{lemma}
\label{lem.partner_partner}
\lean{AlgebraicCombinatorics.SignedCounting.partner_partner}
\leanhelper
\leanok
The partner operation is an involution:
$(I')' = I$ for all finite sets $I$.
\end{lemma}

\begin{proof}
\leanok
By the self-cancellation property of symmetric difference:
$(I \bigtriangleup \{0\}) \bigtriangleup \{0\} = I$.
\end{proof}

\begin{lemma}
\label{lem.partner_sign}
\lean{AlgebraicCombinatorics.SignedCounting.partner_sign}
\leanhelper
\leanok
The partner reverses the sign:
$\operatorname{sign}(I') = -\operatorname{sign}(I)$ for all finite
sets $I$.
\end{lemma}

\begin{proof}
\leanok
If $0 \in I$, then $|I'| = |I| - 1$; if $0 \notin I$, then
$|I'| = |I| + 1$. In either case,
$(-1)^{|I'|} = -(-1)^{|I|}$.
\end{proof}

\begin{proposition}
[Negative hockey-stick identity]
\label{prop.binom.nhs}
\lean{AlgebraicCombinatorics.SignedCounting.negHockeyStick}
\leantarget
\leanok
Let $n\in\mathbb{C}$ and
$m\in\mathbb{N}$. Then,
\[
\sum_{k=0}^{m} (-1)^{k} \binom{n}{k}
= (-1)^{m} \binom{n-1}{m}.
\]
\end{proposition}

\begin{proof}
The Lean proof delegates to
\texttt{Int.alternating\_sum\_range\_choose\_eq\_choose} from Mathlib.

The source gives a combinatorial proof using sign-reversing involutions.
Both sides of the equality are polynomial functions in $n$, so we
can WLOG assume that $n$ is a positive integer (by the polynomial
identity trick). Set $[n] := \{1, 2, \ldots, n\}$.
Define an \emph{acceptable set} to be a subset of $[n]$ that has
size $\leq m$. Define the \emph{sign} of a finite set $I$ to be
$(-1)^{|I|}$. Then
\[
\bigl(\text{the sum of the signs of all acceptable sets}\bigr)
= \sum_{k=0}^{m} (-1)^k \binom{n}{k}.
\]

For each finite set $I$, define the \emph{partner} of $I$ to be the set
$I' := I \bigtriangleup \{1\}$,
where $\bigtriangleup$ denotes symmetric difference.
Each finite set $I$ satisfies $I'' = I$ and
$|I'| = |I| \pm 1$, so that
$(-1)^{|I'|} = -(-1)^{|I|}$. If both $I$ and $I'$ are acceptable
sets, then their contributions to the sum cancel each other.

\textit{Claim 1:} Let $I$ be an acceptable set. Then, the partner
$I'$ of $I$ is non-acceptable if and only if
$(1 \notin I \text{ and } |I| = m)$.

[\textit{Proof of Claim 1:} If $1 \notin I$ and $|I| = m$, then
$I' = I \cup \{1\}$ has size $|I| + 1 > m$, so it is non-acceptable.
Conversely, if $I'$ is non-acceptable, then $1 \notin I$ (otherwise
$I' = I \setminus \{1\} \subseteq I$ would be acceptable) and
$|I| = m$ (otherwise $|I| \leq m - 1$ would give
$|I'| = |I| + 1 \leq m$, making $I'$ acceptable).]

The acceptable sets with non-acceptable partners are precisely the
$m$-element subsets of $[n]$ that do not contain $1$, i.e., the
$m$-element subsets of $[n] \setminus \{1\}$. There are
$\binom{n-1}{m}$ such subsets, each with sign $(-1)^m$.
Hence,
\[
\sum_{k=0}^{m} (-1)^{k} \binom{n}{k} = (-1)^{m} \binom{n-1}{m}.
\]
This proves Proposition~\ref{prop.binom.nhs}.
\end{proof}

\begin{theorem}
\label{thm.acceptable_nonAcceptablePartner_iff}
\lean{AlgebraicCombinatorics.SignedCounting.acceptable_nonAcceptablePartner_iff}
\leanhelper
\leanok
Let $n \geq 1$, $m \in \mathbb{N}$, and let $I$ be an acceptable set
(i.e., $I \subseteq \{0, \ldots, n-1\}$ with $|I| \leq m$). Then the
partner $I'$ of $I$ is non-acceptable if and only if $0 \notin I$
and $|I| = m$.
\end{theorem}

\begin{proof}
\leanok
By case analysis on whether $0 \in I$ and the
cardinality of $I$, as outlined in Claim~1 in the proof of
Proposition~\ref{prop.binom.nhs}.
\end{proof}

\begin{theorem}
\label{thm.sum_signs_acceptable}
\lean{AlgebraicCombinatorics.SignedCounting.sum_signs_acceptable}
\leanhelper
\leanok
For $n \geq 1$ and $m \in \mathbb{N}$,
\[
\sum_{I \in \mathcal{A}} \operatorname{sign}(I)
= (-1)^m \binom{n-1}{m},
\]
where $\mathcal{A} = \mathcal{A}(n, m)$ is the set of acceptable sets
(subsets of $\{0, \ldots, n-1\}$ with size $\leq m$) and
$\operatorname{sign}(I) = (-1)^{|I|}$.
\end{theorem}

\begin{proof}
Partition $\mathcal{A}$ into ``moving'' sets (those whose partner is
also acceptable) and ``fixed'' sets (those whose partner is
non-acceptable). The sum over moving sets is $0$ by a sign-reversing
involution argument: the partner pairs them up with opposite signs.
The sum over fixed sets equals $(-1)^m \binom{n-1}{m}$, since by
Theorem~\ref{thm.acceptable_nonAcceptablePartner_iff}, the fixed sets
are precisely the $m$-element subsets of $[n] \setminus \{0\}$, each
contributing $(-1)^m$.
\end{proof}

\subsection{Cancellation principles}

\begin{lemma}
[Cancellation principle, take 1]
\label{lem.sign.cancel1}
\lean{AlgebraicCombinatorics.SignedCounting.sign_cancel1}
\leantarget
\leanok
Let $\mathcal{A}$ be a finite set. Let $\mathcal{X}$ be a subset of
$\mathcal{A}$.

For each $I \in \mathcal{A}$, let $\operatorname{sign} I$ be an element
of an additive abelian group $R$ with no $2$-torsion (i.e., $2a = 0$
implies $a = 0$). Let $f : \mathcal{X} \to \mathcal{X}$ be a bijection
with the property that
\[
\operatorname{sign}(f(I)) = -\operatorname{sign} I
\qquad \text{for all } I \in \mathcal{X}.
\]
(Such a bijection $f$ is called \emph{sign-reversing}.) Then,
\[
\sum_{I \in \mathcal{A}} \operatorname{sign} I
= \sum_{I \in \mathcal{A} \setminus \mathcal{X}} \operatorname{sign} I.
\]
\end{lemma}

\begin{proof}
We have
\begin{align*}
\sum_{I \in \mathcal{X}} \operatorname{sign} I
&= \sum_{I \in \mathcal{X}} \underbrace{\operatorname{sign}(f(I))}_{
= -\operatorname{sign} I}
\qquad \text{(substituting $f(I)$ for $I$, since $f$ is a bijection)} \\
&= \sum_{I \in \mathcal{X}} (-\operatorname{sign} I)
= -\sum_{I \in \mathcal{X}} \operatorname{sign} I.
\end{align*}
Adding $\sum_{I \in \mathcal{X}} \operatorname{sign} I$ to both sides,
we obtain $2 \cdot \sum_{I \in \mathcal{X}} \operatorname{sign} I = 0$.
Since $R$ has no $2$-torsion, $\sum_{I \in \mathcal{X}} \operatorname{sign} I = 0$.

Now, $\mathcal{X} \subseteq \mathcal{A}$; hence,
\[
\sum_{I \in \mathcal{A}} \operatorname{sign} I
= \underbrace{\sum_{I \in \mathcal{X}} \operatorname{sign} I}_{= 0}
+ \sum_{I \in \mathcal{A} \setminus \mathcal{X}} \operatorname{sign} I
= \sum_{I \in \mathcal{A} \setminus \mathcal{X}} \operatorname{sign} I.
\]
\end{proof}

\begin{lemma}
[Cancellation principle, take 2]
\label{lem.sign.cancel2}
\lean{AlgebraicCombinatorics.SignedCounting.sign_cancel2}
\leantarget
\leanok
Let $\mathcal{A}$ be a finite set. Let $\mathcal{X}$ be a subset of
$\mathcal{A}$.

For each $I \in \mathcal{A}$, let $\operatorname{sign} I$ be an element
of some additive abelian group. Let
$f : \mathcal{X} \to \mathcal{X}$ be an involution (i.e., a map
satisfying $f \circ f = \operatorname{id}$) that has no fixed points.
Assume that
\[
\operatorname{sign}(f(I)) = -\operatorname{sign} I
\qquad \text{for all } I \in \mathcal{X}.
\]
Then,
\[
\sum_{I \in \mathcal{A}} \operatorname{sign} I
= \sum_{I \in \mathcal{A} \setminus \mathcal{X}} \operatorname{sign} I.
\]
\end{lemma}

\begin{proof}
All addends corresponding to the $I \in \mathcal{X}$ cancel out from
the sum $\sum_{I \in \mathcal{A}} \operatorname{sign} I$ (because they
come in pairs of addends with opposite signs). The Lean proof uses
Mathlib's \texttt{Finset.sum\_involution}.
\end{proof}

\begin{lemma}
[Cancellation principle, take 3]
\label{lem.sign.cancel3}
\lean{AlgebraicCombinatorics.SignedCounting.sign_cancel3}
\leantarget
\leanok
Let $\mathcal{A}$ be a finite set. Let $\mathcal{X}$ be a subset of
$\mathcal{A}$.

For each $I \in \mathcal{A}$, let $\operatorname{sign} I$ be an element
of some additive abelian group. Let
$f : \mathcal{X} \to \mathcal{X}$ be an involution (i.e., a map
satisfying $f \circ f = \operatorname{id}$). Assume that
\[
\operatorname{sign}(f(I)) = -\operatorname{sign} I
\qquad \text{for all } I \in \mathcal{X}.
\]
Assume furthermore that
\[
\operatorname{sign} I = 0
\qquad \text{for all } I \in \mathcal{X} \text{ satisfying } f(I) = I.
\]
Then,
\[
\sum_{I \in \mathcal{A}} \operatorname{sign} I
= \sum_{I \in \mathcal{A} \setminus \mathcal{X}} \operatorname{sign} I.
\]
\end{lemma}

\begin{proof}
This is similar to Lemma~\ref{lem.sign.cancel2}, except that the
addends corresponding to the $I \in \mathcal{X}$ satisfying $f(I) = I$
don't cancel (but are already zero and thus can be removed right away).
The Lean proof uses Mathlib's \texttt{Finset.sum\_involution}, which
handles both the pairing and the fixed-point condition.
\end{proof}

\subsection{Switching operations and the \texorpdfstring{$q$}{q}-binomial at \texorpdfstring{$q = -1$}{q = -1}}

\begin{definition}
\label{def.switch}
\lean{AlgebraicCombinatorics.SignedCounting.switch}
\leanhelper
\leanok
For $i \in \mathbb{N}$, define the \emph{switch operation}
$\operatorname{switch}_i$ on finite subsets $S$ of $\mathbb{N}$ by
\[
\operatorname{switch}_i(S) :=
\begin{cases}
S \bigtriangleup \{i, i+1\}, & \text{if } |S \cap \{i, i+1\}| = 1; \\
S, & \text{otherwise.}
\end{cases}
\]
In other words, if exactly one of $i$ and $i+1$ belongs to $S$,
then we replace it by the other; otherwise we leave $S$ unchanged.
\end{definition}

\begin{lemma}
\label{lem.switch_involution}
\lean{AlgebraicCombinatorics.SignedCounting.switch_switch}
\leanhelper
\leanok
The switch operation $\operatorname{switch}_i$ is an involution, i.e.,
$\operatorname{switch}_i(\operatorname{switch}_i(S)) = S$ for all
finite sets $S$.
\end{lemma}

\begin{proof}
\leanok
By case analysis: if exactly one of $\{i, i+1\}$ is in $S$, then the
symmetric difference $S \bigtriangleup \{i, i+1\}$ restores $S$ after
two applications. Otherwise the operation is the identity.
\end{proof}

\begin{lemma}
\label{lem.switch_card}
\lean{AlgebraicCombinatorics.SignedCounting.switch_card}
\leanhelper
\leanok
The switch operation preserves cardinality:
$|\operatorname{switch}_i(S)| = |S|$ for all finite sets $S$.
\end{lemma}

\begin{proof}
\leanok
When $|S \cap \{i, i+1\}| = 1$, the symmetric difference removes one
element and adds another. Otherwise the set is unchanged.
\end{proof}

\subsection{Roots of unity}

\begin{definition}
\label{def.root-of-unity.prim}
\lean{AlgebraicCombinatorics.SignedCounting.IsRootOfUnity, IsPrimitiveRoot}
\leantarget
\leanok
Let $K$ be a field. Let $d$ be a positive integer.

\textbf{(a)} A \emph{$d$-th root of unity} in $K$ means an element
$\omega$ of $K$ satisfying $\omega^d = 1$.

\textbf{(b)} A \emph{primitive $d$-th root of unity} in $K$ means an
element $\omega$ of $K$ satisfying $\omega^d = 1$ but
$\omega^i \neq 1$ for each $i \in \{1, 2, \ldots, d-1\}$.
In other words, a primitive $d$-th root of unity in $K$ means an element
of the multiplicative group $K^{\times}$ whose order is $d$.
\end{definition}

\begin{lemma}
\label{lem.neg_one_isPrimitiveRoot_two}
\lean{AlgebraicCombinatorics.SignedCounting.neg_one_isPrimitiveRoot_two}
\leanhelper
\leanok
In any integral domain where $-1 \neq 1$, the element $-1$ is a
primitive $2$nd root of unity.
\end{lemma}

\begin{proof}
\leanok
We have $(-1)^2 = 1$ and $(-1)^1 = -1 \neq 1$.
\end{proof}

\begin{lemma}
\label{lem.primitiveRoot_geom_sum}
\lean{AlgebraicCombinatorics.SignedCounting.primitiveRoot_geom_sum_eq_zero}
\leanhelper
\leanok
If $\omega$ is a primitive $d$-th root of unity with $d > 1$, then
\[
1 + \omega + \omega^2 + \cdots + \omega^{d-1} = 0.
\]
\end{lemma}

\begin{proof}
\leanok
This is \texttt{IsPrimitiveRoot.geom\_sum\_eq\_zero} from Mathlib.
\end{proof}

\subsection{The \texorpdfstring{$q$}{q}-Lucas theorem}

\begin{definition}
\label{def.qBinomial}
\lean{AlgebraicCombinatorics.SignedCounting.qBinomial}
\leanhelper
\leanok
The \emph{$q$-binomial coefficient} $\binom{n}{k}_q$ is the polynomial
in $\mathbb{Z}[q]$ defined by
\[
\binom{n}{k}_q := \sum_{\substack{S \subseteq \{0,\ldots,n-1\} \\ |S| = k}}
q^{\,\operatorname{sum}(S) - (0 + 1 + \cdots + (k-1))},
\]
where $\operatorname{sum}(S) = \sum_{x \in S} x$.
This is the local definition used in the formalization; it agrees with
Mathlib's $q$-binomial coefficient (see
Lemma~\ref{lem.qBinomial_eq_canonical}).
\end{definition}

\begin{lemma}
\label{lem.qBinomial_eq_canonical}
\lean{AlgebraicCombinatorics.SignedCounting.qBinomial_eq_canonical}
\leanhelper
\leanok
The local $q$-binomial coefficient
$\binom{n}{k}_q$ (Definition~\ref{def.qBinomial}) equals
Mathlib's $q$-binomial coefficient $\binom{n}{k}_q$ (as a polynomial
in $\mathbb{Z}[q]$).
\end{lemma}

\begin{proof}
\leanok
Both definitions agree: they are sums over $k$-element subsets of
$\{0,\ldots,n-1\}$, with monomials
$q^{(\operatorname{sum}(S) - k(k-1)/2)}$.
The proof establishes the identity by showing the exponents coincide
using the order embedding of $S$ into $\{0,\ldots,n-1\}$.
\end{proof}

\begin{lemma}
\label{lem.qBinomial_eval_one}
\lean{AlgebraicCombinatorics.SignedCounting.qBinomial_eval_one}
\leanhelper
\leanok
For $n, k \in \mathbb{N}$,
$\binom{n}{k}_q \big|_{q=1} = \binom{n}{k}$.
\end{lemma}

\begin{proof}
\leanok
Each term evaluates to $1$ at $q = 1$, so the sum counts the number of
$k$-element subsets of $[n]$, which is $\binom{n}{k}$.
\end{proof}

\begin{lemma}
\label{lem.qBinomial_eq_zero_of_lt}
\lean{AlgebraicCombinatorics.SignedCounting.qBinomial_eq_zero_of_lt}
\leanhelper
\leanok
For $n < k$, $\binom{n}{k}_q = 0$.
\end{lemma}

\begin{proof}
\leanok
There are no $k$-element subsets of $[n]$ when $n < k$.
\end{proof}

\begin{lemma}
\label{lem.qBinomial_zero_right}
\lean{AlgebraicCombinatorics.SignedCounting.qBinomial_zero_right}
\leanhelper
\leanok
For all $n \in \mathbb{N}$, $\binom{n}{0}_q = 1$.
\end{lemma}

\begin{proof}
\leanok
The only $0$-element subset is $\varnothing$, contributing
$q^0 = 1$.
\end{proof}

\begin{lemma}
\label{lem.qBinomial_self}
\lean{AlgebraicCombinatorics.SignedCounting.qBinomial_self}
\leanhelper
\leanok
For all $n \in \mathbb{N}$, $\binom{n}{n}_q = 1$.
\end{lemma}

\begin{proof}
\leanok
The only $n$-element subset of $[n]$ is $[n]$ itself, contributing
$q^0 = 1$.
\end{proof}

\begin{definition}
\label{def.IsDBlocky}
\lean{AlgebraicCombinatorics.SignedCounting.IsDBlocky}
\leanhelper
\leanok
A subset $S \subseteq \{0, \ldots, n-1\}$ is \emph{$d$-blocky} if
$S \subseteq \{0, \ldots, n-1\}$ and for each block index
$i \in \{0, \ldots, \lfloor n/d \rfloor - 1\}$, either all elements of
$\{id, id+1, \ldots, id+d-1\}$ are in $S$ or none are.
\end{definition}

\begin{definition}
\label{def.kSubsets}
\lean{AlgebraicCombinatorics.SignedCounting.kSubsets}
\leanhelper
\leanok
The set $\binom{[n]}{k}$ of $k$-element subsets of
$\{0, \ldots, n-1\}$:
\[
\binom{[n]}{k} := \bigl\{S \subseteq \{0,\ldots,n-1\} : |S| = k\bigr\}.
\]
\end{definition}

\begin{definition}
\label{def.blockySubsets}
\lean{AlgebraicCombinatorics.SignedCounting.blockySubsets}
\leanhelper
\leanok
The set of $d$-blocky $k$-element subsets of $\{0, \ldots, n-1\}$:
\[
\mathcal{B}(d, n, k) :=
\bigl\{S \in \binom{[n]}{k} : S \text{ is } d\text{-blocky}\bigr\}.
\]
\end{definition}

\begin{lemma}
\label{lem.blocky_subsets_count}
\lean{AlgebraicCombinatorics.SignedCounting.blocky_subsets_count}
\leanhelper
\leanok
For $d > 0$, the number of $d$-blocky $k$-element subsets of
$\{0, \ldots, n-1\}$ is
\[
|\mathcal{B}(d, n, k)|
= \begin{cases}
\binom{n/d}{k/d} \cdot \binom{n \bmod d}{k \bmod d},
  & \text{if } k \bmod d \leq n \bmod d; \\
0, & \text{otherwise.}
\end{cases}
\]
\end{lemma}

\begin{proof}
\leanok
A $d$-blocky $k$-element subset is uniquely determined by the pair
$(B, P)$ where $B \subseteq \{0, \ldots, n/d - 1\}$ selects which
complete $d$-blocks are fully included (giving $|B| \cdot d$ elements)
and $P \subseteq \{0, \ldots, n \bmod d - 1\}$ selects which elements
of the partial block are included (giving $|P|$ elements).
The constraint $|S| = k$ forces $|B| = k / d$ and $|P| = k \bmod d$.
If $k \bmod d > n \bmod d$, no valid $P$ exists.
\end{proof}

\begin{theorem}
[$q$-Lucas theorem]
\label{thm.sign.q-lucas}
\lean{AlgebraicCombinatorics.SignedCounting.qLucas}
\leantarget
\leanok
Let $K$ be a field. Let $d$ be a positive integer. Let $\omega$ be a
primitive $d$-th root of unity in $K$. Let $n, k \in \mathbb{N}$. Let
$q$ and $u$ be the quotients obtained when dividing $n$ and $k$ by $d$
with remainder, and let $r$ and $v$ be the respective remainders. Then,
\[
\binom{n}{k}_{\omega}
= \binom{q}{u} \cdot \binom{r}{v}_{\omega}.
\]
\end{theorem}

\begin{proof}
The proof splits into cases:

\textbf{Case $d = 1$:} Then $\omega = 1$, and both sides reduce to
$\binom{n}{k}$ by Lemma~\ref{lem.qBinomial_eval_one}.

\textbf{Case $k > n$:} Both sides are $0$.

\textbf{Case $n < d$:} Then $n / d = 0$, $n \bmod d = n$, and the
formula reduces to $\binom{n}{k}_\omega = \binom{0}{k/d} \cdot
\binom{n}{k \bmod d}_\omega$.

\textbf{Main case ($d \geq 2$, $k \leq n$, $n \geq d$):}
The proof has three steps:
\begin{enumerate}
\item Rewrite $\binom{n}{k}_\omega$ as $\sum_{S \in \mathcal{A}}
\omega^{(\operatorname{sum}(S) - \text{baseline})}$ where $\mathcal{A}$
is the set of $k$-element subsets of $[n]$.
\item Split the sum into contributions from \emph{$d$-blocky} subsets
(those that consist of entire $d$-blocks) and non-$d$-blocky subsets.
\item Show that non-blocky contributions cancel
(Lemma~\ref{lem.nonBlocky_contributions_cancel}) and that the blocky
contributions give the product formula
(Lemma~\ref{lem.blocky_sum_eq}).
\end{enumerate}
A subset $S$ of $[n]$ is \emph{$d$-blocky} if for each complete
$d$-block $\{id, id+1, \ldots, id+d-1\}$ contained in $[n]$, either
all elements of the block are in $S$ or none are.
\end{proof}

\begin{lemma}
\label{lem.nonBlocky_contributions_cancel}
\lean{AlgebraicCombinatorics.SignedCounting.nonBlocky_contributions_cancel}
\leanhelper
\leanok
Let $K$ be a field, $d \geq 2$, and $\omega$ a primitive $d$-th root
of unity. For $n, k \in \mathbb{N}$,
\[
\sum_{\substack{S \subseteq [n],\; |S| = k,\\ S \text{ not $d$-blocky}}}
\omega^{\operatorname{sum}(S)} = 0.
\]
\end{lemma}

\begin{proof}
For each non-$d$-blocky set $S$, there exists a ``split block''
$i$: a block index where $S$ contains some but not all elements
of $\{id, id+1, \ldots, id+d-1\}$. Define the rotation operation
that cyclically permutes elements within block $i$ (replacing each
element $x$ in the block by $id + (x - id + 1) \bmod d$).

The non-blocky sets are partitioned into orbits under rotation in
the smallest split block. Each orbit has size
$m = d / \gcd(d, j)$ where $j = |S \cap \text{block}_i|$, and
$m > 1$ since $0 < j < d$. The sum over each orbit is $0$ by
Lemma~\ref{lem.orbit_sum_eq_zero}. Therefore the total sum is $0$.

The Lean proof applies Lemma~\ref{lem.sum_nonBlocky_subset_eq_zero}
to the full set of non-blocky subsets, verifying orbit-closure
(rotation preserves membership via
Lemmas~\ref{lem.rotateSetInBlockK_card} and
\ref{lem.rotateSetInBlockK_subset}).
\end{proof}

\begin{lemma}
\label{lem.orbit_sum_eq_zero}
Let $K$ be a field, $d \geq 2$, and $\omega$ a primitive $d$-th root
of unity. For $0 < j < d$ and any $b \in \mathbb{N}$, let
$m = d / \gcd(d, j)$. Then
\[
\sum_{k=0}^{m-1} \omega^{b + k \cdot j} = 0.
\]
\end{lemma}

\begin{proof}
Factor out $\omega^b$:
$\sum_{k=0}^{m-1} \omega^{b + kj} = \omega^b \sum_{k=0}^{m-1} (\omega^j)^k$.
Since $0 < j < d$, we have $\omega^j \neq 1$ and $(\omega^j)^m = 1$.
The sum $\sum_{k=0}^{m-1} (\omega^j)^k$ can be evaluated using the
geometric series formula
$\frac{(\omega^j)^m - 1}{\omega^j - 1} = 0$.
\end{proof}

\begin{lemma}
\label{lem.blocky_sum_eq}
\lean{AlgebraicCombinatorics.SignedCounting.blocky_sum_eq}
\leanhelper
\leanok
Let $K$ be a field, $d$ a positive integer, and $\omega$ a primitive
$d$-th root of unity. For $n, k \in \mathbb{N}$,
\[
\sum_{\substack{S \subseteq [n],\; |S| = k,\\ S \text{ is $d$-blocky}}}
\omega^{\operatorname{sum}(S) - (0 + 1 + \cdots + (k-1))}
= \binom{n/d}{k/d} \cdot
\binom{n \bmod d}{k \bmod d}_{\omega}.
\]
\end{lemma}

\begin{proof}
A $d$-blocky $k$-element subset of $[n]$ decomposes into two
independent parts: which complete $d$-blocks to include (choosing
$k/d$ blocks from $n/d$ available) and which elements of the partial
block $\{(n/d) \cdot d, \ldots, n-1\}$ to include (choosing
$k \bmod d$ elements from $n \bmod d$ available).

The complete-block part contributes a factor of $1$ (since
$\omega^d = 1$) and the partial-block part contributes
$\omega^{(\text{sum of partial part})}$. The number of
complete-block configurations is $\binom{n/d}{k/d}$, and the
partial-block sum gives
$\binom{n \bmod d}{k \bmod d}_{\omega}$.
\end{proof}

\begin{lemma}
\label{lem.rotateInBlock_injective}
The rotation function within block $i$ is injective (for $d > 0$).
\end{lemma}

\begin{proof}
By case analysis on whether $x$ and $y$ are in block $i$, and using
injectivity of the modular arithmetic operation.
\end{proof}

\begin{lemma}
\label{lem.rotateInBlock_iterate_d}
Rotating $d$ times within block $i$ returns to the identity:
$(\operatorname{rotate}_i)^d(x) = x$ for all $x$.
\end{lemma}

\begin{proof}
For $x$ in block $i$, after $k$ rotations the offset becomes
$(x - id + k) \bmod d$. After $d$ rotations, the offset returns to
its original value. For $x$ outside block $i$, the rotation is the
identity.
\end{proof}

\begin{lemma}
\label{lem.rotateSetInBlockK_sum_pow}
For a primitive $d$-th root of unity $\omega$,
\[
\omega^{\operatorname{sum}(\operatorname{rotate}_i^k(S))}
= \omega^{\operatorname{sum}(S) + k \cdot |\text{blockOffsets}(d, i, S)|}.
\]
\end{lemma}

\begin{proof}
The rotation changes $\operatorname{sum}(S)$ by
$|\text{blockOffsets}(d, i, S)| \pmod{d}$ per iteration. Since
$\omega^d = 1$, the congruence modulo $d$ suffices.
\end{proof}

\begin{lemma}
\label{lem.rotateSetInBlockK_injective}
For $0 < j < d$ where $j = |\text{blockOffsets}(d, i, S)|$ and
$m = d / \gcd(d, j)$, the map
$k \mapsto \operatorname{rotate}_i^k(S)$ is injective on
$\{0, 1, \ldots, m-1\}$.
\end{lemma}

\begin{proof}
If $\operatorname{rotate}_i^{k_1}(S) = \operatorname{rotate}_i^{k_2}(S)$
with $k_1 < k_2 < m$, then $k_1 \cdot j \equiv k_2 \cdot j \pmod{d}$.
By modular arithmetic cancellation, $m \mid (k_2 - k_1)$, contradicting
$0 < k_2 - k_1 < m$.
\end{proof}

\begin{lemma}
\label{lem.rotateSetInBlockK_card}
Rotation preserves cardinality: for $d > 0$,
$|\operatorname{rotate}_i^k(S)| = |S|$.
\end{lemma}

\begin{proof}
The rotation function $\operatorname{rotate}_i$ is injective
(Lemma~\ref{lem.rotateInBlock_injective}), so the image of $S$
under $\operatorname{rotate}_i$ has the same cardinality.
The result follows by induction on $k$.
\end{proof}

\begin{lemma}
\label{lem.rotateSetInBlockK_subset}
Rotation preserves the ambient set: for $d > 0$ and $i < n/d$,
if $S \subseteq \{0, \ldots, n-1\}$ then
$\operatorname{rotate}_i^k(S) \subseteq \{0, \ldots, n-1\}$.
\end{lemma}

\begin{proof}
Elements outside block $i$ are unchanged. Elements inside block $i$
stay within $\{id, \ldots, id+d-1\} \subseteq \{0, \ldots, n-1\}$
by the modular arithmetic of the rotation.
\end{proof}

\begin{lemma}
\label{lem.rotateSetInBlockK_d}
Rotating $d$ times is the identity on sets:
$\operatorname{rotate}_i^d(S) = S$ for all sets $S$.
\end{lemma}

\begin{proof}
Since $(\operatorname{rotate}_i)^d = \operatorname{id}$ on elements
(Lemma~\ref{lem.rotateInBlock_iterate_d}), the image of $S$ after $d$
rotations equals $S$.
\end{proof}

\begin{lemma}
\label{lem.orbit_sum_full_eq_zero}
Let $\omega$ be a primitive $d$-th root of unity with $d \geq 2$.
For $0 < j < d$ and any $b \in \mathbb{N}$,
\[
\sum_{k=0}^{d-1} \omega^{b + k \cdot j} = 0.
\]
\end{lemma}

\begin{proof}
Since $m = d / \gcd(d, j)$ divides $d$, the sum over the full range
$\{0, \ldots, d-1\}$ is a sum of $d/m$ copies of the orbit sum
$\sum_{k=0}^{m-1} \omega^{b + kj}$, which is $0$ by
Lemma~\ref{lem.orbit_sum_eq_zero}.
\end{proof}

\begin{lemma}
\label{lem.rotateSetInBlockK_fiber_card_uniform}
For $d > 0$ and any set $S$, the fibers of the map
$k \mapsto \operatorname{rotate}_i^k(S)$ on $\{0, \ldots, d-1\}$ all
have the same cardinality. That is, there exists $c$ such that for
every set $T$ in the image, the number of $k \in \{0, \ldots, d-1\}$
with $\operatorname{rotate}_i^k(S) = T$ is exactly $c$, and $c \mid d$.
\end{lemma}

\begin{proof}
By the orbit-stabilizer theorem for the cyclic group $\mathbb{Z}/d$
acting on sets via rotation. The stabilizer of $S$ consists of those
$k$ with $\operatorname{rotate}_i^k(S) = S$; all fibers have the same
size as the stabilizer. The result follows from
$|\operatorname{orbit}| \cdot |\operatorname{stab}| = d$.
\end{proof}

\begin{lemma}
\label{lem.sum_nonBlocky_subset_eq_zero}
For any orbit-closed subset $A$ of the non-blocky $k$-element subsets
of $[n]$,
\[
\sum_{S \in A} \omega^{\operatorname{sum}(S)} = 0.
\]
Here ``orbit-closed'' means that $A$ is closed under rotation in the
smallest split block of each element: for each $S \in A$, the entire
orbit of $S$ under rotation in its smallest split block is contained
in $A$.
\end{lemma}

\begin{proof}
By strong induction on $|A|$. If $A$ is empty, the sum is $0$.
Otherwise, pick $S_0 \in A$, let $i_0$ be its smallest split block
index, and let $O$ be the orbit of $S_0$ under rotation in block
$i_0$ (all $d$ rotations). Since $A$ is orbit-closed, $O \subseteq A$.
The sum over $O$ is $0$ by
Lemma~\ref{lem.orbit_sum_full_eq_zero} and
Lemma~\ref{lem.rotateSetInBlockK_fiber_card_uniform}
(the sum over the image equals a constant times the sum over the
fibers, which is zero). Then
$\sum_{S \in A} \omega^{\operatorname{sum}(S)} = \sum_{S \in O}(\cdots) + \sum_{S \in A \setminus O}(\cdots) = 0 + 0$
by the inductive hypothesis on $A \setminus O$ (which is orbit-closed
and has strictly smaller cardinality).
\end{proof}

\subsection{The \texorpdfstring{$q$}{q}-binomial at \texorpdfstring{$q = -1$}{q = -1}}

\begin{theorem}
\label{thm.qBinom_neg_one}
\lean{AlgebraicCombinatorics.SignedCounting.qBinom_neg_one}
\leanhelper
\leanok
Let $n, k \in \mathbb{N}$. Then:
\[
\binom{n}{k}_{-1} =
\begin{cases}
0, & \text{if } n \text{ is even and } k \text{ is odd;} \\
\binom{\lfloor n/2 \rfloor}{\lfloor k/2 \rfloor}, & \text{otherwise.}
\end{cases}
\]
\end{theorem}

\begin{proof}
By Theorem~\ref{thm.sign.q-lucas} with $d = 2$ and $\omega = -1$
(which is a primitive $2$nd root of unity by
Lemma~\ref{lem.neg_one_isPrimitiveRoot_two}):
\[
\binom{n}{k}_{-1}
= \binom{n/2}{k/2} \cdot \binom{n \bmod 2}{k \bmod 2}_{-1}.
\]
The only possible remainders upon division by $2$ are $0$ and $1$, and
the base-case $(-1)$-binomial coefficients are:
\[
\binom{0}{0}_{-1} = 1, \quad
\binom{0}{1}_{-1} = 0, \quad
\binom{1}{0}_{-1} = 1, \quad
\binom{1}{1}_{-1} = 1.
\]
If $n$ is even and $k$ is odd, then $n \bmod 2 = 0$ and
$k \bmod 2 = 1$, so $\binom{n \bmod 2}{k \bmod 2}_{-1} = 0$ and
the result is $0$. Otherwise,
$\binom{n \bmod 2}{k \bmod 2}_{-1} = 1$ and the result is
$\binom{\lfloor n/2 \rfloor}{\lfloor k/2 \rfloor}$.

The Lean proof applies the $q$-Lucas theorem (Theorem~\ref{thm.sign.q-lucas})
with $d = 2$ over $\mathbb{Q}$, then uses injectivity of
$\mathbb{Z} \hookrightarrow \mathbb{Q}$ to transfer the result back
to $\mathbb{Z}$.
\end{proof}
